\chapter{Data Structures at Scale: How Industry Builds on Classical Foundations}

\section{Introduction}

Classical data structures---red-black trees, hash tables, B-trees, skip lists---form
the theoretical backbone of every software system.  A textbook treatment typically
assumes that data fits in memory and that operations run on a single machine.  In
practice, companies like Google, Amazon, Facebook, and Netflix operate systems that
serve billions of requests per second across thousands of nodes.  At that scale,
the gap between a clean $O(\log n)$ analysis and a production-grade system is
enormous.

Two forces drive the gap.  First, hardware is not uniform: SSDs favour sequential I/O,
DRAM is scarce relative to data size, and multi-core NUMA machines make cache
locality a first-order concern.  Second, distribution introduces partial failure:
nodes crash, networks partition, and consistency guarantees must be negotiated
explicitly.  Classical data structures were not designed with these constraints in
mind.

This chapter bridges the gap.  We study six systems---LSM trees, consistent hashing,
column-family storage, Bloom filter variants, copy-on-write semantics, and
hardware-aware design---that illustrate how industry adapts classical ideas.  Each
section includes a C++ code sketch, a complexity discussion, and pointers to the
papers and open-source projects that made these ideas real.  By the end, you will
have a practical vocabulary for choosing and tuning data structures when the dataset
no longer fits on one machine or in one disk.

\section{Log-Structured Merge Trees --- The Engine Behind Modern Databases}

\subsection{Motivation}

Write-heavy workloads---event logging, time-series ingestion, IoT telemetry---stress
traditional B-tree databases because every insertion may trigger a random disk write
to maintain sorted order.  Log-Structured Merge (LSM) trees trade read performance
for write throughput by batching writes in memory and flushing them sequentially to
disk.

\subsection{Architecture}

An LSM tree consists of three components:
\begin{enumerate}
    \item \textbf{Memtable:} an in-memory sorted structure (typically a red-black
    tree or skip list) that absorbs all writes.
    \item \textbf{SSTables (Sorted String Tables):} immutable, sorted files on disk,
    each containing a sequence of key-value pairs together with an index block and
    optional Bloom filters.
    \item \textbf{Compaction process:} a background thread that merges adjacent
    SSTables to reclaim space and reduce the number of files that reads must examine.
\end{enumerate}

\subsection{Write Path}

Every write enters the memtable first.  Once the memtable reaches a size threshold
(typically 4--64\,MB), it is flushed to disk as a new SSTable.  The flush is
sequential, so even on spinning disks it is fast.  The amortized write cost is
$O(1)$ for the memtable insertion plus the amortized cost of background compaction.

\begin{lstlisting}[language=C++]
#include <map>
#include <vector>
#include <string>
#include <fstream>
#include <algorithm>

struct SSTable {
    std::vector<std::pair<std::string, std::string>> data;

    void writeToFile(const std::string& path) const {
        std::ofstream out(path);
        for (const auto& [key, value] : data) {
            out << key << "\t" << value << "\n";
        }
    }

    static SSTable merge(const SSTable& a, const SSTable& b) {
        SSTable result;
        std::merge(a.data.begin(), a.data.end(),
                   b.data.begin(), b.data.end(),
                   std::back_inserter(result.data));
        return result;
    }
};

class LSMTree {
    std::map<std::string, std::string> memtable;
    std::vector<SSTable> levels[4];
    size_t memtableLimit = 1000;

public:
    void put(const std::string& key, const std::string& value) {
        memtable[key] = value;
        if (memtable.size() >= memtableLimit) {
            flushMemtable();
        }
    }

    std::string get(const std::string& key) {
        if (memtable.count(key)) return memtable[key];

        for (int i = 0; i < 4; ++i) {
            for (int j = levels[i].size() - 1; j >= 0; --j) {
                auto& table = levels[i][j];
                auto it = std::lower_bound(
                    table.data.begin(), table.data.end(),
                    std::make_pair(key, std::string()),
                    [](const auto& a, const auto& b) {
                        return a.first < b.first;
                    });
                if (it != table.data.end() && it->first == key) {
                    return it->second;
                }
            }
        }
        return "";
    }

private:
    void flushMemtable() {
        SSTable sst;
        for (const auto& [k, v] : memtable) {
            sst.data.emplace_back(k, v);
        }
        levels[0].push_back(sst);
        memtable.clear();
        compact(0);
    }

    void compact(int level) {
        if (levels[level].size() < 4) return;
        SSTable merged = levels[level][0];
        for (size_t i = 1; i < levels[level].size(); ++i) {
            merged = SSTable::merge(merged, levels[level][i]);
        }
        levels[level].clear();
        levels[level + 1].push_back(merged);
    }
};
\end{lstlisting}

\subsection{Read Path}

A read checks the memtable first, then SSTables from newest to oldest.  Without
auxiliary structures this is $O(k \cdot \log m)$ where $k$ is the number of SSTables
and $m$ is the size of each.  Bloom filters attached to each SSTable reduce the
expected number of disk probes to close to $1$ for point queries.

\subsection{Compaction Strategies}

Two dominant strategies exist.  \emph{Size-tiered compaction} (used by Cassandra
and HBase) groups SSTables of similar size and merges them when a threshold is
reached.  \emph{Leveled compaction} (used by LevelDB and RocksDB) organises SSTables
into levels with a size ratio of 10:1 and merges individual SSTables from level $i$
into level $i+1$.  Levelled compaction uses more I/O but produces fewer tombstones
and gives more predictable read performance.

\subsection{Real-World Implementations}

Google's Bigtable (2006) was the first production LSM tree at scale.  LevelDB
(2011) made the design accessible as an embeddable library, and RocksDB (2013,
Facebook) extended it with compression, transactions, and column families.  As of
2024, RocksDB is the storage engine behind MySQL (MyRocks), Kafka Streams, and
CockroachDB.

\subsection{LSM vs.\ B-tree: A Comparison}

B-trees provide $O(\log n)$ reads with a single disk seek because keys live at a
known location in a single sorted tree.  LSM trees defer sorting to compaction, so
writes are sequential and batching reduces the total I/O volume.  On SSDs, the
sequential-write advantage is smaller but still meaningful because SSD garbage
collection prefers contiguous data.  LSM trees generally win when the
write-to-read ratio exceeds 10:1; B-trees win for read-heavy or latency-sensitive
workloads.

\section{Consistent Hashing --- Distributing Data Without Central Control}

\subsection{The Problem}

Suppose you have $N$ cache servers and you distribute keys with $h(k) \bmod N$.
When a server is added or removed, $N$ changes and nearly every key maps to a
different server, causing a thundering-restore of cached data.

\subsection{The Ring}

Consistent hashing (Karger et al., 1997) hashes both keys and server identifiers
onto a ring with positions $[0, 2^{32})$.  A key is assigned to the next server
clockwise from its hash position.  When a server joins, only the keys in the arc
between it and its predecessor need to be remapped.

\begin{lstlisting}[language=C++]
#include <map>
#include <string>
#include <vector>
#include <functional>
#include <algorithm>
#include <numeric>

class ConsistentHashRing {
    std::map<size_t, std::string> ring;
    std::hash<std::string> hasher;
    int virtualNodes;

public:
    explicit ConsistentHashRing(int vnodes = 150)
        : virtualNodes(vnodes) {}

    void addNode(const std::string& node) {
        for (int i = 0; i < virtualNodes; ++i) {
            std::string vkey = node + "#" + std::to_string(i);
            ring[hasher(vkey)] = node;
        }
    }

    void removeNode(const std::string& node) {
        for (int i = 0; i < virtualNodes; ++i) {
            std::string vkey = node + "#" + std::to_string(i);
            ring.erase(hasher(vkey));
        }
    }

    std::string getNode(const std::string& key) const {
        size_t h = hasher(key);
        auto it = ring.lower_bound(h);
        if (it == ring.end()) it = ring.begin();
        return it->second;
    }

    std::vector<std::string> getReplicas(
            const std::string& key, int count) const {
        std::vector<std::string> replicas;
        size_t h = hasher(key);
        auto it = ring.lower_bound(h);
        std::set<std::string> seen;
        while (replicas.size() < count) {
            if (it == ring.end()) it = ring.begin();
            if (seen.insert(it->second).second) {
                replicas.push_back(it->second);
            }
            ++it;
        }
        return replicas;
    }
};
\end{lstlisting}

\subsection{Virtual Nodes}

With only one hash per physical server, the arc lengths on the ring follow an
exponential distribution, causing severe imbalance.  Virtual nodes (vnodes) assign
each physical server multiple positions on the ring.  With 150 vnodes per server,
the standard deviation of load drops below 5\% in practice.

\subsection{Production Use}

Amazon's DynamoDB paper (DeCandia et al., 2007) introduced consistent hashing to
a wide audience.  Today the technique underpins Cassandra's partitioner, Memcached's
ketama algorithm, Netflix's EVCache, and Uber's Ringpop library.

\section{Column-Family Storage --- Wide Columns at Scale}

\subsection{Data Model}

A column-family store generalises the relational model.  Data is indexed by a row
key, then a column family, then a column qualifier, and optionally a timestamp:

\[
    \text{Row Key} \to \text{Column Family} \to \text{Column Qualifier}
    \to \text{Timestamp} \to \text{Value}
\]

This model is \emph{sparse}: rows need not store every column, and storage cost is
proportional to the data actually written, not to a schema-wide column count.

\subsection{Bigtable and Its Descendants}

Google's Bigtable paper (Chang et al., 2006) described a distributed storage system
built on GFS that used this model.  HBase (Apache, 2008) brought the design to
Hadoop, and Cassandra (Facebook, 2008) combined the wide-column model with
consistent hashing for peer-to-peer replication.

\begin{lstlisting}[language=C++]
#include <map>
#include <string>
#include <vector>
#include <cstdint>

struct Cell {
    std::string columnQualifier;
    int64_t timestamp;
    std::string value;
};

using ColumnFamily = std::map<std::string, std::vector<Cell>>;
using RowKey = std::string;

class ColumnFamilyStore {
    std::map<RowKey,
        std::map<std::string, std::vector<Cell>>> table;

public:
    void put(const RowKey& row, const std::string& cf,
             const std::string& col, const std::string& val) {
        int64_t ts = currentTimestamp();
        table[row][cf].push_back({col, ts, val});
    }

    std::string get(const RowKey& row, const std::string& cf,
                    const std::string& col) const {
        auto rowIt = table.find(row);
        if (rowIt == table.end()) return "";

        auto cfIt = rowIt->second.find(cf);
        if (cfIt == rowIt->second.end()) return "";

        const auto& cells = cfIt->second;
        for (auto it = cells.rbegin(); it != cells.rend(); ++it) {
            if (it->columnQualifier == col) return it->value;
        }
        return "";
    }

    std::vector<std::pair<std::string, std::string>>
    scan(const RowKey& row, const std::string& cf) const {
        std::vector<std::pair<std::string, std::string>> result;
        auto rowIt = table.find(row);
        if (rowIt == table.end()) return result;

        auto cfIt = rowIt->second.find(cf);
        if (cfIt == rowIt->second.end()) return result;

        for (const auto& cell : cfIt->second) {
            result.emplace_back(cell.columnQualifier, cell.value);
        }
        return result;
    }

private:
    int64_t currentTimestamp() const {
        return std::chrono::steady_clock::now().time_since_epoch()
            .count();
    }
};
\end{lstlisting}

\subsection{Why Column Families Win for Analytics}

Column families enable \emph{column pruning}: a query over three columns reads only
those three columns, not entire rows.  Adjacent values in a column family compress
well because they share data types and value ranges.  For time-series workloads,
a typical query reads one row key (a sensor ID) and a single column family (the
metric), achieving $O(\log n)$ lookup followed by sequential reads within the column
family.

\section{Bloom Filters in Production}

\subsection{Recap}

A Bloom filter is a probabilistic set that answers \emph{``Is $x$ in the set?''}
with no false negatives and a tunable false-positive rate $\varepsilon$.  Given $m$
bits and $k$ hash functions, $\varepsilon \approx (1 - e^{-kn/m})^k$.

\subsection{At Scale}

Google uses Bloom filters throughout its storage stack.  Bigtable attaches a Bloom
filter to each SSTable to avoid reading files that provably do not contain a queried
key.  Reported reductions in disk reads range from 100x to 1000x for random lookups.
Netflix uses Bloom filters to deduplicate content recommendations across its
recommendation engine, reducing redundant API calls.

\subsection{Alternatives}

Cuckoo filters support deletion (impossible with standard Bloom filters) and offer
better cache locality because they store fingerprints in buckets rather than
scattering them across a bit array.  Quotient filters achieve similar benefits with
a compact, cache-friendly layout.  Both have become popular in production systems
where the data set evolves and items must be removed.

\section{Copy-on-Write and Versioned Storage}

\subsection{Persistent Data Structures via CoW}

Copy-on-write (CoW) allows a structure to be shared among multiple readers while
each writer sees its own version.  The key insight is to copy only the path from the
root to the modified node, reusing the rest.

Git is the most widely deployed CoW data structure.  Each commit object points to a
tree of blobs; unchanged subtrees are shared between commits.  File systems like ZFS
and Btrfs use CoW at the block level to provide atomic snapshots and crash
consistency.

\begin{lstlisting}[language=C++]
#include <memory>
#include <vector>
#include <string>

template <typename T>
class CowPtr {
    std::shared_ptr<const T> data;

public:
    explicit CowPtr(std::shared_ptr<const T> d)
        : data(std::move(d)) {}

    const T& read() const { return *data; }

    CowPtr clone() const {
        return CowPtr(std::make_shared<T>(*data));
    }

    CowPtr write(auto&& modifier) const {
        auto copy = std::make_shared<T>(*data);
        modifier(*copy);
        return CowPtr(std::move(copy));
    }
};

struct Config {
    std::string dbPath;
    int maxConnections;
    bool enableCompression;

    Config() : dbPath("/data/db"),
               maxConnections(64),
               enableCompression(true) {}
};

void versionedConfig() {
    CowPtr<Config> v1(std::make_shared<Config>());
    CowPtr<Config> v2 = v1.write([](Config& c) {
        c.maxConnections = 128;
    });
    CowPtr<Config> v3 = v2.write([](Config& c) {
        c.enableCompression = false;
    });
    std::cout << v1.read().maxConnections << "\n";
    std::cout << v3.read().enableCompression << "\n";
}
\end{lstlisting}

\subsection{Production Systems}

Cassandra's commit log and memtable together form a shadow-copy system: the commit
log is append-only, and the memtable is an in-memory snapshot that can be replaced
atomically.  This design gives Cassandra write-ahead-log durability without the
write amplification of a B-tree.

\section{Priority Queues and Scheduling at Scale}

\subsection{Kafka and Message Priorities}

Apache Kafka traditionally orders messages by offset within a partition.  When
priority ordering is needed, operators maintain separate topics per priority level
and use a priority consumer that drains higher-priority partitions first.  At scale,
this is effectively a distributed priority queue implemented with partition
assignment.

\subsection{Redis Sorted Sets}

Redis sorted sets store a member with an associated score and use a skip list
combined with a hash table internally.  The skip list provides $O(\log n)$ insert,
delete, and rank queries, while the hash table gives $O(1)$ membership lookups.
Redis's in-memory nature makes this a natural choice for leaderboards, job queues,
and scheduling in game servers.

\begin{lstlisting}[language=C++]
#include <queue>
#include <vector>
#include <string>
#include <fstream>
#include <functional>

struct Job {
    int priority;
    std::string taskId;
    bool operator<(const Job& o) const {
        return priority < o.priority;
    }
};

class DiskSpillPQ {
    std::priority_queue<Job> inMemory;
    std::vector<std::string> spillFiles;
    size_t memoryLimit = 10000;
    size_t currentSize = 0;
    int fileCounter = 0;

public:
    void push(const Job& job) {
        if (currentSize >= memoryLimit) {
            spillToDisk();
        }
        inMemory.push(job);
        ++currentSize;
    }

    bool pop(Job& result) {
        if (inMemory.empty()) {
            if (spillFiles.empty()) return false;
            loadFromDisk();
        }
        result = inMemory.top();
        inMemory.pop();
        --currentSize;
        return true;
    }

private:
    void spillToDisk() {
        std::string path = "/tmp/pq_spill_" +
            std::to_string(fileCounter++) + ".bin";
        std::ofstream out(path, std::ios::binary);
        while (!inMemory.empty()) {
            Job j = inMemory.top();
            inMemory.pop();
            out.write(reinterpret_cast<const char*>(&j.priority),
                      sizeof(int));
            size_t len = j.taskId.size();
            out.write(reinterpret_cast<const char*>(&len),
                      sizeof(size_t));
            out.write(j.taskId.data(), len);
        }
        spillFiles.push_back(path);
        currentSize = 0;
    }

    void loadFromDisk() {
        std::string path = spillFiles.back();
        spillFiles.pop_back();
        std::ifstream in(path, std::ios::binary);
        while (in.peek() != EOF) {
            Job j;
            in.read(reinterpret_cast<char*>(&j.priority),
                    sizeof(int));
            size_t len;
            in.read(reinterpret_cast<char*>(&len), sizeof(size_t));
            j.taskId.resize(len);
            in.read(j.taskId.data(), len);
            inMemory.push(j);
            ++currentSize;
        }
        std::remove(path.c_str());
    }
};
\end{lstlisting}

\section{Graph Data Structures in Social Networks}

\subsection{The Social Graph}

Facebook's TAO (The Associations and Objects) paper (Bronson et al., 2013) describes
a distributed graph store serving over a billion users.  TAO stores a directed
multigraph where vertices represent users, pages, and posts, and edges represent
likes, follows, and comments.  Queries are expressed as graph traversals with
pagination.

\subsection{Property Graph Model}

A property graph associates typed properties with both vertices and edges.  This
model underpins Neo4j, Amazon Neptune, and TigerGraph.  The key design question is
how to store edges efficiently: as adjacency lists (good for traversals from one
vertex) or as compressed sparse row/column (CSR/CSC) formats (good for bulk
analysis and memory density).

\begin{lstlisting}[language=C++]
#include <vector>
#include <string>
#include <unordered_map>
#include <algorithm>

struct Edge {
    int target;
    std::string label;
};

class CompressedGraph {
    std::vector<int> rowPtr;
    std::vector<int> colIdx;
    std::vector<Edge> edges;
    std::vector<std::string> vertexProps;
    int numVertices;

public:
    explicit CompressedGraph(int n)
        : numVertices(n), rowPtr(n + 1, 0) {}

    void addEdge(int from, int to,
                 const std::string& label) {
        edges.push_back({to, label});
        rowPtr[from + 1]++;
    }

    void buildIndex() {
        for (int i = 1; i <= numVertices; ++i) {
            rowPtr[i] += rowPtr[i - 1];
        }
        std::vector<int> offsets = rowPtr;
        std::vector<Edge> sorted(edges.size());
        for (const auto& e : edges) {
            int pos = offsets[0]++;
            sorted[pos] = e;
        }
        edges = std::move(sorted);
        colIdx.resize(edges.size());
        for (size_t i = 0; i < edges.size(); ++i) {
            colIdx[i] = edges[i].target;
        }
    }

    std::vector<int> neighbors(int vertex) const {
        int start = rowPtr[vertex];
        int end = rowPtr[vertex + 1];
        return std::vector<int>(
            colIdx.begin() + start,
            colIdx.begin() + end);
    }

    int degree(int vertex) const {
        return rowPtr[vertex + 1] - rowPtr[vertex];
    }
};
\end{lstlisting}

\subsection{Graph Partitioning}

Distributing a graph across machines requires partitioning.  Hash-based partitioning
is simple ($\text{partition}(v) = h(v) \bmod k$) but cuts many edges, increasing
cross-machine traversals.  METIS and power-law-aware algorithms cluster high-degree
vertices near the graph centre, reducing the edge cut at the expense of some load
imbalance.  TAO uses a two-level scheme: vertex IDs encode the shard, and edges are
co-located with their source vertex.

\section{Hardware-Aware Data Structure Design}

\subsection{Cache Line Awareness}

A modern x86 cache line is 64 bytes.  Two threads writing to different fields of
the same struct that share a cache line cause \emph{false sharing}, degrading
performance by orders of magnitude.  Aligning frequently written fields to cache
line boundaries eliminates this:

\begin{lstlisting}[language=C++]
#include <new>
#include <cstdint>
#include <cstdlib>

struct alignas(64) CacheLinePadded {
    int64_t counter;
    char padding[64 - sizeof(int64_t)];
};

struct ThreadLocalCounters {
    CacheLinePadded perThread[64];
    int64_t sum() const {
        int64_t total = 0;
        for (int i = 0; i < 64; ++i) {
            total += perThread[i].counter;
        }
        return total;
    }
};

void* allocateAligned(size_t alignment, size_t size) {
    void* ptr = nullptr;
    posix_memalign(&ptr, alignment, size);
    return ptr;
}

struct alignas(64) GraphNode {
    int vertexId;
    int64_t distance;
    int parentIndex;
    int edgeStart;
    int edgeCount;
    char pad[64 - 5 * sizeof(int)];
};
\end{lstlisting}

\subsection{NUMA-Aware Allocation}

On multi-socket systems, memory attached to a remote socket costs 2--3x more in
latency.  The first-touch policy assigns pages to the NUMA node of the thread that
first writes to them.  Spreading initialisation across threads on all sockets (or
using \texttt{mbind} with interleave policy) avoids hot-spotting a single socket.

\subsection{SOA vs AOS}

Array-of-Structs (AOS) stores all fields of each element contiguously.
Structure-of-Arrays (SOA) stores each field in a separate contiguous array.
For SIMD vectorisation or queries that touch one field, SOA avoids loading
irrelevant data into vector registers.  A graph's adjacency list, for instance,
benefits from SOA when only the target vertex IDs are needed for a breadth-first
search.

\section{Comparison Table and Design Decisions}

\begin{center}
\begin{tabular}{@{} l l l l @{}}
\toprule
\textbf{Structure} & \textbf{Write} & \textbf{Read} & \textbf{Best For} \\
\midrule
LSM Tree & $O(1)$ amort. & $O(\log n)$ with Bloom & Write-heavy, time-series \\
B-Tree & $O(\log n)$ & $O(\log n)$ single-seek & Read-heavy, transactions \\
Hash Index & $O(1)$ amort. & $O(1)$ expected & Point lookups, caches \\
Column-Family & $O(\log n)$ & $O(\log n + k)$ & Sparse, wide, analytics \\
\bottomrule
\end{tabular}
\end{center}

\begin{center}
\begin{tabular}{@{} l l l l @{}}
\toprule
\textbf{Model} & \textbf{Schema} & \textbf{Scaling} & \textbf{Use Case} \\
\midrule
Row Store & Fixed, rigid & Vertical / sharding & OLTP, transactions \\
Column Family & Semi-structured & Horizontal, automatic & IoT, time-series \\
Document Store & Schema-free & Horizontal, automatic & CMS, user profiles \\
\bottomrule
\end{tabular}
\end{center}

\begin{center}
\textbf{Decision Flowchart for Data Structure Selection at Scale}
\end{center}

Start with the access pattern.  If the workload is write-dominant
(write:read $> 10$:1) and point queries suffice, an LSM tree is the default choice.
If reads dominate and latency is critical, prefer a B-tree.  If the data is
schema-free and document-shaped, a document store is natural.  For wide, sparse
time-series data, column families with compaction tuned for your read/write ratio
are optimal.  If the data is a graph with dense traversals, a native graph store
outperforms relational joins.  Finally, if keys must be distributed across an
elastic set of servers, always use consistent hashing with virtual nodes rather than
modulo partitioning.

\section{Exercises}

\subsection{Drill}

\begin{enumerate}
    \item An LSM tree has 3 levels with a size ratio of 10.  If the memtable is
    4\,MB and each level holds $10^i \times 4$\,MB, what is the total on-disk
    capacity before compaction to the next level?  How many SSTables are examined
    per point query without Bloom filters?

    \item A consistent hash ring has 10 physical nodes, each with 150 virtual
    nodes.  What is the expected number of virtual nodes on the ring?  If a new
    node is added, approximately what fraction of keys are remapped?

    \item A Bloom filter has 10 million bits and 7 hash functions for a set of
    1 million elements.  Calculate the false-positive probability.  How many bits
    per element is this?
\end{enumerate}

\subsection{Application}

\begin{enumerate}
    \item Implement the LSM tree from Section~23.2 and benchmark insertion of
    10 million random key-value pairs with and without Bloom filters on SSTables.
    Measure both write throughput and point-query latency.  Plot the results as a
    function of the number of SSTables.

    \item Build a column-family store using the API from Section~23.4.  Ingest a
    synthetic time-series dataset (1000 sensors, 86400 readings each).  Compare
    the scan time of reading all readings for one sensor against reading all
    sensors for one timestamp.  Explain the performance difference in terms of
    column-family layout.

    \item Implement the disk-spilling priority queue from Section~23.7 with a
    threshold of 1000 in-memory jobs.  Push 1 million jobs and then pop them all.
    Measure the time spent spilling and reloading compared to a pure in-memory
    priority queue.
\end{enumerate}

\subsection{Research}

\begin{enumerate}
    \item Read the RocksDB Wiki section on compaction strategies.  Implement a
    simple size-tiered compaction simulator that models write amplification, space
    amplification, and read amplification as functions of the number of levels and
    the size ratio.  Compare your model's predictions against empirical benchmarks
    published in the RocksDB documentation.

    \item The Facebook TAO paper describes a two-level caching architecture with
    a global cache backed by MySQL.  Read the paper (Bronson et al., 2013) and
    sketch an LRU cache implementation in C++ that supports batched reads to
    minimise round-trip overhead.  Analyse the amortised cost per lookup under
    varying cache-hit rates.
\end{enumerate}

%% ============================================================
%% INDUSTRIAL PERSPECTIVES (merged from former Chapter 24)
%% ============================================================

\section{Industrial Perspectives}

Every data structure and algorithm in this book was chosen because it appears in real systems. But not all are equally important in practice. The table below provides an honest assessment of each topic's industrial relevance, rated on a three-point scale:

\begin{center}
\begin{tabular}{cl}
\toprule
Rating & Meaning \\
\midrule
\textbf{3 -- High} & Used daily in production systems at scale \\
\textbf{2 -- Medium} & Used in specific domains or as building blocks \\
\textbf{1 -- Low} & Primarily educational, theoretical, or superseded \\
\bottomrule
\end{tabular}
\end{center}

These ratings reflect the state of industry in 2025. A rating of ``Low'' does not mean a topic is unimportant for learning---it means you are unlikely to encounter it in production code.

%% --- Master Relevance Table ---

\subsection{Complete Relevance Table}

\begin{center}
\small
\begin{tabular}{lcl}
\toprule
\textbf{Data Structure / Algorithm} & \textbf{Rating} & \textbf{Primary Industrial Use} \\
\midrule
\multicolumn{3}{l}{\textit{Linear Structures}} \\
\midrule
Arrays / Dynamic Arrays & 3 & Everywhere (std::vector, database pages) \\
Linked Lists & 2 & Kernel internals, LRU caches \\
Stacks & 3 & Compilers, DFS, undo/redo, call stacks \\
Queues & 3 & BFS, message queues, task scheduling \\
Deques & 3 & Sliding window, BFS, schedulers \\
Circular Buffers & 3 & Ring buffers, networking, producer-consumer \\
Gap Buffers & 1 & Text editors (Emacs, vi) \\
Unrolled Lists & 1 & Some database systems \\
Finger Trees & 1 & Functional programming (Haskell) \\
Ropes & 1 & Text editors, some languages \\
\midrule
\multicolumn{3}{l}{\textit{Hashing}} \\
\midrule
Hash Tables (Chaining) & 3 & Ubiquitous (databases, caches, compilers) \\
Hash Tables (Open Addressing) & 3 & High-performance systems (Google, Redis) \\
Cuckoo Hashing & 2 & Concurrent hash tables, network routers \\
Robin Hood Hashing & 2 & Some STL implementations \\
Swiss Table & 3 & Google Abseil, high-performance C++ \\
Skip Lists & 3 & Redis sorted sets, concurrent data structures \\
Count-Min Sketch & 2 & Stream processing, network monitoring \\
Merkle Trees (Hash Trees) & 3 & Blockchain, distributed systems, Git \\
Extendible Hashing & 2 & Database file organizations \\
Distributed Hash Table & 2 & P2P systems, CDNs \\
\midrule
\multicolumn{3}{l}{\textit{Search Trees}} \\
\midrule
Binary Search Trees & 3 & Foundation for all balanced trees \\
AVL Trees & 3 & Database indexing, some std::set \\
Red-Black Trees & 3 & Java TreeMap, C++ std::map, Linux CFS \\
B-Trees & 3 & Databases (PostgreSQL, MySQL), filesystems \\
B+-Trees & 3 & Database indexing, NTFS, ext4 \\
B*-Trees & 2 & Some database systems \\
Scapegoat Trees & 1 & Educational \\
Treaps & 2 & Concurrent data structures, some databases \\
Splay Trees & 1 & Some caching systems, mostly educational \\
Van Emde Boas Trees & 1 & Theoretical, some integer-specific systems \\
Weight-Balanced Trees & 1 & Some database systems \\
(a,b)-Trees & 2 & B-tree variants in databases \\
Top-Down Red-Black & 2 & Implementation variant \\
Finger Trees & 1 & Functional programming niche \\
Partial Rebuilding & 1 & Educational \\
Optimal BST & 1 & Educational \\
\midrule
\multicolumn{3}{l}{\textit{Spatial Trees}} \\
\midrule
K-d Trees & 3 & Graphics, ML, nearest neighbor search \\
Quadtrees & 2 & 2D spatial indexing, games \\
Octrees & 2 & 3D graphics, game engines \\
R-Trees & 3 & GIS, database spatial indexing (PostGIS) \\
BSP Trees & 1 & Game rendering, mostly replaced \\
Range Trees & 2 & Computational geometry \\
\midrule
\multicolumn{3}{l}{\textit{Graph Structures}} \\
\midrule
Adjacency Matrix & 3 & Dense graphs, Floyd-Warshall \\
Adjacency List & 3 & Sparse graphs, most applications \\
Hypergraphs & 1 & Specialized (circuit design) \\
BDDs & 1 & Formal verification, circuit design \\
Winged Edge & 1 & Mesh representation, CAD \\
\midrule
\multicolumn{3}{l}{\textit{Graph Algorithms}} \\
\midrule
BFS & 3 & Shortest paths, web crawling, social networks \\
DFS & 3 & Topological sort, SCC, cycle detection \\
Dijkstra's Algorithm & 3 & GPS, routing, networks \\
Bellman-Ford & 2 & Routing with negative weights \\
Floyd-Warshall & 2 & All-pairs, small graphs \\
Prim's MST & 3 & Network design, clustering \\
Kruskal's MST & 3 & Network design, clustering \\
Topological Sort & 3 & Build systems, task scheduling, compilers \\
Strongly Connected Components & 3 & Web graph analysis, compilers \\
A* Search & 3 & Games, robotics, GPS navigation \\
Graph Coloring & 2 & Scheduling, register allocation \\
Eulerian Path & 1 & Postal routing, puzzles \\
Articulation Points / Bridges & 2 & Network reliability \\
\midrule
\multicolumn{3}{l}{\textit{String Algorithms}} \\
\midrule
Tries & 3 & Autocomplete, IP routing, spell checking \\
KMP & 2 & Text search (replaced by suffix arrays) \\
Rabin-Karp & 2 & Plagiarism detection, multi-pattern search \\
Boyer-Moore & 2 & Text editors (grep) \\
Suffix Arrays & 3 & Search engines, bioinformatics \\
Suffix Trees & 2 & Replaced by suffix arrays in most cases \\
Aho-Corasick & 2 & IDS, antivirus, plagiarism detection \\
Z-Algorithm & 1 & Replaced by suffix arrays \\
Document Retrieval & 3 & Search engines (Google, Bing) \\
DC3/Skew Algorithm & 2 & Suffix array construction \\
\midrule
\multicolumn{3}{l}{\textit{Range Query Structures}} \\
\midrule
Sparse Table & 2 & Static RMQ \\
Segment Tree & 3 & Databases, competitive programming, range queries \\
Fenwick Tree & 3 & Inversion counting, range sums \\
Interval Tree & 2 & Scheduling, interval overlap \\
Square Root Decomposition & 2 & Competitive programming \\
2D Fenwick Tree & 2 & Image processing \\
Orthogonal Range Trees & 2 & Computational geometry \\
Segment Tree Beats & 1 & Competitive programming \\
\midrule
\multicolumn{3}{l}{\textit{Sorting Algorithms}} \\
\midrule
Insertion Sort & 3 & Small arrays, nearly sorted, std::sort fallback \\
Selection Sort & 1 & Educational \\
Bubble Sort & 1 & Educational \\
Merge Sort & 3 & Stable sort, external sort, Java Arrays.sort \\
Quick Sort & 3 & std::sort, most common general-purpose sort \\
Heap Sort & 2 & Guaranteed O(n log n), introsort fallback \\
Radix Sort & 3 & Database engines, string sorting \\
Counting Sort & 2 & Integer sorting \\
Bucket Sort & 2 & Uniform distribution sorting \\
\midrule
\multicolumn{3}{l}{\textit{Greedy Algorithms}} \\
\midrule
Huffman Coding & 3 & Compression (ZIP, JPEG, MP3, HTTP/2) \\
Activity Selection & 2 & Scheduling \\
Fractional Knapsack & 1 & Educational \\
Job Sequencing & 2 & Task scheduling \\
\midrule
\multicolumn{3}{l}{\textit{Dynamic Programming}} \\
\midrule
0/1 Knapsack & 2 & Resource allocation \\
Matrix Chain Multiplication & 2 & Database query optimization \\
LCS & 3 & Diff tools, bioinformatics, version control \\
Edit Distance & 3 & Spell checking, DNA alignment, fuzzy search \\
Coin Change & 2 & Financial systems \\
LIS & 2 & Patience sorting, scheduling \\
Floyd-Warshall (DP) & 2 & All-pairs shortest paths \\
\midrule
\multicolumn{3}{l}{\textit{Probabilistic Structures}} \\
\midrule
Bloom Filter & 3 & Databases, CDNs, networking (Google, Meta) \\
HyperLogLog & 3 & Analytics (Redis, PostgreSQL, BigQuery) \\
MinHash & 2 & Document similarity, deduplication \\
LSH & 2 & Similarity search, recommendation systems \\
Consistent Hashing & 3 & Distributed systems, CDNs, DynamoDB \\
\midrule
\multicolumn{3}{l}{\textit{Network Flow}} \\
\midrule
Max Flow / Min Cut & 2 & Transportation, scheduling, image segmentation \\
Bipartite Matching & 2 & Assignment problems, scheduling \\
Min-Cost Max-Flow & 2 & Logistics, supply chain \\
\midrule
\bottomrule
\end{tabular}
\end{center}

%% --- Tier Descriptions ---

\subsection{Tier 3 --- High Relevance: The Industrial Workhorses}

These data structures and algorithms are used daily in production systems at companies like Google, Meta, Amazon, Microsoft, Netflix, and virtually every technology company.

\textbf{Hash tables} are the single most important data structure in industry. They power database index lookups (MySQL, PostgreSQL, MongoDB), caching (Redis, Memcached, CDN edge caches), compiler symbol tables, web server session management, network router forwarding tables, and operating system process tables. Google's Abseil library uses Swiss Table hash maps, which are 2--3$\times$ faster than \texttt{std::unordered\_map}.

\textbf{B-trees and B+-trees} are the backbone of every major database and filesystem. PostgreSQL uses B+-trees as the default index type; MySQL InnoDB's clustered index is a B+-tree; ext4 and NTFS use B-trees for directory indexing. A single PostgreSQL query can traverse millions of B+-tree nodes per second.

\textbf{Red-black trees} are the default balanced tree in most standard libraries: Java's \texttt{TreeMap}, C++'s \texttt{std::map}, and Linux's Completely Fair Scheduler. The Linux kernel's CFS scheduler traverses a red-black tree of millions of processes millions of times per second.

\textbf{Segment trees and Fenwick trees} power real-time analytics and range queries in database engines, competitive programming, image processing, and time-series databases.

\textbf{Graph algorithms} (BFS, DFS, Dijkstra, MST, topological sort, SCC) are essential infrastructure. Google Maps uses a variant of Dijkstra's on a graph with 100 million+ intersections.

\textbf{Tries} power Google's autocomplete (3.5 billion queries/day), IP routing, and spell checking. \textbf{Bloom filters} are used by Chrome (malicious URL detection), Cassandra (SSTable lookup), and Bitcoin (SPV verification). \textbf{Consistent hashing} is the foundation of DynamoDB, Cassandra, CDNs, and Memcached.

\subsection{Tier 2 --- Medium Relevance: Domain Specialists}

\textbf{Skip lists} are the backbone of Redis sorted sets and LevelDB/RocksDB memtables. \textbf{Suffix arrays} have replaced suffix trees for text indexing and bioinformatics (BWA aligns 1 billion DNA reads in 2 hours). \textbf{Max flow} solves transportation, image segmentation, and airline scheduling problems. \textbf{HyperLogLog} counts 1 billion distinct elements in 12~KB (Redis, PostgreSQL, BigQuery). \textbf{Persistent data structures} underpin Git, undo systems, and MVCC in databases. \textbf{LSM trees} power LevelDB, RocksDB, Cassandra, and InfluxDB---RocksDB processes 10 million writes/second on a single SSD.

\subsection{Tier 1 --- Low Relevance: Educational and Theoretical}

\textbf{Splay trees} have good amortized performance but are mostly superseded by red-black trees. \textbf{Van Emde Boas trees} achieve $O(\lg \lg u)$ but hash tables are faster in practice. \textbf{Fusion trees} prove a theoretical lower bound is achievable but have no practical implementation. \textbf{Selection sort and bubble sort} are purely educational. The ``low relevance'' structures are not wasted knowledge---understanding them builds deep theoretical intuition that makes you a better engineer.

%% --- Decision Matrix ---

\subsection{What to Learn for Industry}

\begin{center}
\small
\begin{tabular}{ll}
\toprule
\textbf{Career Path} & \textbf{Must-Know Structures} \\
\midrule
Backend Engineering & Hash tables, B-trees, sorting, BFS/DFS, Dijkstra \\
Database Engineering & B+-trees, LSM trees, hash indexes, external sort \\
Search Engineering & Tries, suffix arrays, inverted indices, Bloom filters \\
Game Development & K-d trees, A*, quaternions, spatial hashing \\
Machine Learning & Sparse matrices, hash tables, k-d trees, LSH \\
Distributed Systems & Consistent hashing, Bloom filters, vector clocks \\
Embedded Systems & Circular buffers, bit arrays, succinct structures \\
Competitive Programming & Segment trees, Fenwick trees, union-find, DP \\
\bottomrule
\end{tabular}
\end{center}

\subsection{Industrial Perspectives---Exercises}

\begin{enumerate}
\item \textbf{Industry survey.} Pick a large open-source project (e.g., Linux kernel, Redis, PostgreSQL, Chromium). Identify 10 data structures it uses. Classify each by relevance rating.

\item \textbf{Startup stack.} You are building a real-time analytics platform processing 1 million events/second. Which data structures from this book would you choose for: (a) counting unique visitors, (b) maintaining top-$k$ trending topics, (c) storing user profiles with fast lookup, (d) detecting duplicate events?

\item \textbf{Database internals.} Explain why B+-trees are preferred over red-black trees for database indexing, even though red-black trees have better worst-case bounds. What role does caching play?

\item \textbf{Distributed design.} You are designing a distributed cache with 100 nodes. Explain why consistent hashing is essential and what happens without it when a node fails.

\item \textbf{Trade-off analysis.} For each tier-2 structure, identify one system where it is irreplaceable and one system where a tier-1 alternative could work. Justify your choices.
\end{enumerate}

\section{References and Further Reading}

\begin{itemize}
    \item Chang, F., Dean, J., Ghemawat, S., et al.\ (2006).  Bigtable: A
    Distributed Storage System for Structured Data.  \emph{Proc.\ OSDI}.

    \item DeCandia, G., Hastorun, D., Jampani, M., et al.\ (2007).  Dynamo:
    Amazon's Highly Available Key-value Store.  \emph{Proc.\ SOSP}.

    \item Bronson, N., Amsden, Z., Cabral, G., et al.\ (2013).  TAO:
    Facebook's Distributed Data Store for the Social Graph.  \emph{Proc.\ ATC}.

    \item O'Neil, P., Cheng, E., Gawlick, D., and O'Neil, E.\ (1996).
    The Log-Structured Merge-Tree (LSM-Tree).  \emph{Acta Informatica}, 33(4),
    351--385.

    \item Karger, D., Sherman, A., Berkheimer, A., et al.\ (1997).  Consistent
    Hashing and Random Trees: Distributed Caching Protocols for Relieving Hot
    Spots on the World Wide Web.  \emph{Proc.\ STOC}.

    \item Kreps, J., Narkhede, N., and Rao, J.\ (2011).  Kafka: A Distributed
    Messaging System for Log Processing.  \emph{Proc.\ NetDB}.

    \item Lakshman, A.\ and Malik, P.\ (2010).  Cassandra: A Decentralized
    Structured Storage System.  \emph{ACM SIGOPS Operating Systems Review}, 44(2),
    35--40.

    \item RocksDB Wiki.  \url{https://github.com/facebook/rocksdb/wiki}

    \item Melnik, S.\ and Gubarev, A.\ (2010).  Dremel: Interactive Analysis of
    Web-Scale Datasets.  \emph{Proc.\ VLDB}, 3, 330--339.

    \item Conchon, S.\ and Filliâtre, J.-C.\ (2008).  A Persistent
    Union-Find Data Structure.  \emph{Proc.\ ML Workshop on ML}, 37--46.

    \item Bender, M. A., Farach-Colton, M., and Kim, B.\ (2019).  Engineering
    a Cache-Oblivious Sort that is Fast in Practice.  \emph{Workshop on Algorithm
    Engineering and Experiments (ALENEX)}.

    \item Dean, J.\ and Barroso, L. A.\ (2013).  The Tail at Scale.
    \emph{Communications of the ACM}, 56(2):74--83.

    \item Bloom, B. H.\ (1970).  Space/Time Trade-offs in Hash Coding with
    Allowable Errors.  \emph{Communications of the ACM}, 13(7):422--426.

    \item Kirsch, A.\ and Mitzenmacher, M.\ (2006).  Less Hashing, Same
    Performance: Building a Better Bloom Filter.  \emph{ACM Transactions on
    Algorithms}, 2(4):589--605.

    \item Leis, V.\ et al.\ (2013).  The Adaptive Radix Tree: Practical Indexing
    for Main-Memory Database Systems.  \emph{IEEE ICDE}.

    \item Stonebraker, M.\ et al.\ (2005).  C-Store: A Column-Oriented DBMS.
    \emph{VLDB}.
\end{itemize}
